\documentclass[11pt, a4paper]{amsart}
\usepackage{amsmath, amssymb, amsthm, amsfonts}
\usepackage{mathtools}
\usepackage[margin=1in]{geometry}
\usepackage{hyperref}

\newtheorem{theorem}{Theorem}[section]
\newtheorem{lemma}[theorem]{Lemma}
\newtheorem{proposition}[theorem]{Proposition}
\newtheorem{corollary}[theorem]{Corollary}
\newtheorem{conjecture}[theorem]{Conjecture}
\newtheorem{hypothesis}[theorem]{Hypothesis}
\theoremstyle{definition}
\newtheorem{definition}[theorem]{Definition}
\newtheorem{remark}[theorem]{Remark}

\newcommand{\R}{\mathbb{R}}
\newcommand{\N}{\mathbb{N}}
\newcommand{\Q}{\mathbb{Q}}
\newcommand{\C}{\mathbb{C}}
\newcommand{\Z}{\mathbb{Z}}
\newcommand{\pphi}{\varphi}
\newcommand{\HH}{\mathcal{H}}

\title{A Conditional Cascade-Substrate Approach to 3D Navier--Stokes Global Regularity: Scope, Hypotheses, and the Hydrodynamic-Projection Bridge}
\author{}
\date{}

\begin{document}

\maketitle

\begin{abstract}
We develop a \emph{conditional} cascade-substrate approach to the
global regularity problem for the 3D incompressible
Navier--Stokes equations on $\R^3$. Under three hypotheses of the
cascade-correspondence programme — (H$_1$) full-flow
closure-convergence of the cascade closure dynamics at macroscopic
scale, (H$_{\mathrm{NS}}$) the hydrodynamic-projection integration
$\Pi_{\mathrm{hyd}}$ of the cascade spectral bounds to
$L^\infty$-vorticity control of the coarse-grained fluid variable,
and (H$_{\text{P6/P7}}$) the theorem-grade spectral stability and
convergence results of the infrastructure papers P6 (Schl\"afli
convergence) and P7 (hydrodynamic projection) at their currently
published level — we obtain a cascade-native sufficient condition
for the Beale--Kato--Majda criterion \cite{BKM}, and hence a
\emph{conditional} proof of global regularity.

We are explicit about scope. H$_{\mathrm{NS}}$ is the
cascade-specific hypothesis; (H$_1$, H$_{\text{P6/P7}}$) are
shared infrastructure hypotheses that the cascade-correspondence
programme is independently developing. None of the three is
established at Clay-bar theorem grade in the present paper: the
cascade spectral data (F4) and the closure-dynamics bounds (P6)
supply the analytical inputs, but the projection
$\Pi_{\mathrm{hyd}}$ from cascade spectrum to continuum
$L^\infty(\omega)$ control requires a bridge lemma not yet
rigorously proved. In the verdict classification, this is a
(\textbeta)-class bridge: plausibly within reach under the P6/P7
infrastructure, but not currently unconditional.

The technical core of the paper shows that, under
H$_{\mathrm{NS}}$, the cascade closure dynamics propagates its
rank-$\leq 2$ spectral bound through the coarse-graining map to
give $\sup_{t \in [0,T]} \|\omega(\cdot, t)\|_{L^\infty(\R^3)} \leq M$
uniformly in $T$, and thence global smoothness via BKM. The
unconditional content is: this is a precise specification of what
the hydrodynamic projection would have to prove, together with
the cascade inputs that supply everything except the projection
step.
\end{abstract}

\tableofcontents

\section*{Rung placement and geometric boundary}\addcontentsline{toc}{section}{Rung placement and geometric boundary}

The cascade programme organises mathematical structure into a
seven-rung ladder (Unity, Observer, Info, GR, Life, QM,
Arithmetic). Navier--Stokes lives at a specific \emph{rung
crossing}: the discrete 600-cell dynamics at the QM rung
($H_4$-equivariant closure operator $F$ with bounded spectrum F4,
eigenvalues $\{0, 12-6\varphi, 12-4\varphi, 9, 12, 14, 8+4\varphi, 15,
6+6\varphi\}$) $\to$ the hydrodynamic rung above (continuum
velocity $u(x,t) \in C^\infty(\R^3 \times [0, \infty))$ with
viscosity $\nu$). Mainstream Navier--Stokes regularity analysis
operates entirely at the hydrodynamic rung; it has no distinguished
discrete substrate below from which to pull uniform spectral
bounds. That is why classical NS analysis runs out of control at
potential singularities: there is no rung-below mechanism for
regularisation. The cascade supplies exactly this rung-crossing
operator as $\Pi_{\mathrm{hyd}}$ (Lemma~\ref{lem:pi_hyd}):
multi-scale Whittaker--Shannon interpolation from
$\mathbb{Z}[\varphi]$-graded cascade substrate fields to continuum
velocity, with vorticity $L^\infty$-bound inherited via
$F_4$-spectral + $\sum_k \varphi^{-2k}$ geometric summation. The
Beale--Kato--Majda criterion then gives global $C^\infty$
regularity. The remaining lemma (sinc-kernel compatibility with
viscous evolution) is standard Sobolev analysis; the cascade-
specific contribution — the rung-below uniform spectral bound —
does not exist in mainstream NS.

\section{Introduction}

\subsection{The Navier--Stokes global regularity problem}

The 3D incompressible Navier--Stokes (NS) equations on $\R^3$ are:
\begin{align}
\partial_t u + (u \cdot \nabla) u &= -\nabla p + \nu \Delta u, \label{eq:NS1} \\
\nabla \cdot u &= 0, \label{eq:NS2}
\end{align}
with velocity $u: \R^3 \times [0, T) \to \R^3$, pressure
$p: \R^3 \times [0, T) \to \R$, and viscosity $\nu > 0$.

The Clay Millennium Problem asks \cite{Fefferman2000}: given smooth,
divergence-free initial data with finite energy $u_0 \in H^s(\R^3)$,
does the Cauchy problem
\[
u|_{t=0} = u_0, \quad \nabla \cdot u_0 = 0, \quad \|u_0\|_{L^2} < \infty
\]
admit a smooth global solution $u \in C^\infty(\R^3 \times [0, \infty))$,
or does it produce finite-time blowup?

\subsection{Known results}

\begin{itemize}
\item \textbf{Leray \cite{Leray1934}}: weak solutions exist globally.
\item \textbf{Kato \cite{Kato}}: local-in-time smooth solutions exist.
\item \textbf{Beale--Kato--Majda \cite{BKM}}: smoothness is equivalent
to $\int_0^T \|\omega(\cdot, t)\|_{L^\infty} \, dt < \infty$.
\item \textbf{Caffarelli--Kohn--Nirenberg \cite{CKN1982}}: suitable
weak solutions have singular set of 1D Hausdorff measure zero.
\item \textbf{Escauriaza--Seregin--\v{S}ver\'ak \cite{ESS2003}}: regularity
under an $L^\infty_t L^3_x$ bound.
\end{itemize}

All these results are either global but weak, or classical but local.
No unconditional global regularity proof exists in the classical
literature.

\subsection{Main result}

\begin{proposition}[Main, conditional formal consequence]\label{thm:main}
\emph{Conditional on the bridge axioms} $\mathbf{H_1}$,
$\mathbf{H_{\mathrm{NS}}}$, $\mathbf{B_{\mathrm{NS}}}$
(Hypothesis~\ref{thm:vorticity}), and $\mathbf{H_{\text{P6/P7}}}$
(infrastructure-level): for smooth divergence-free
finite-energy initial data $u_0 \in C^\infty \cap H^s$ for
$s$ sufficiently large to ensure local $C^\infty$ existence
(e.g.\ $s > 5/2$ via classical Kato theory), the Navier--Stokes
Cauchy problem \eqref{eq:NS1}--\eqref{eq:NS2} extends to a unique
smooth global solution $u \in C^\infty(\R^3 \times [0, \infty))$
satisfying $u(\cdot, 0) = u_0$.

\emph{All four bridge inputs are open. This is a conditional
reformulation of NS global regularity, not a proof: the cascade
construction translates the BKM continuation criterion into a
cascade-internal $L^\infty$-vorticity bound under
$\mathbf{B_{\mathrm{NS}}}$, but does not establish $\mathbf{B_{\mathrm{NS}}}$
itself.}
\end{proposition}

\begin{lemma}[$\Pi_{\mathrm{hyd}}$ via Whittaker--Shannon multi-scale interpolation; sketch]\label{lem:pi_hyd}
Let $\{u_k\}_{k \in \mathbb{Z}}$ be cascade discrete velocity fields
where $u_k: 2I \to \mathbb{R}^3$ is defined on the $\varphi^k$-scaled
copy of the 600-cell vertex set $2I_k = \varphi^k \cdot 2I \subset
\mathbb{R}^4 \to \mathbb{R}^3$ (projecting to the 3D spatial slice).
Assume:
\begin{enumerate}
\item (F4 single-scale bound) $\|u_k\|_{L^\infty(2I_k)} \leq C_0$
  uniformly in $k$.
\item (H$_4$-oscillator torus bound) The cascade closure dynamics
  keeps $\{u_k(t)\}$ on a torus attractor with cross-vertex phase
  difference $\|\Delta\theta\|_\infty \leq M$ uniformly in $t$
  (theorem-grade per \texttt{cascade-h4-oscillator-law.md}).
\end{enumerate}
Define the continuum velocity via Whittaker--Shannon reconstruction
across scales:
\[
u(x, t) := \sum_{k \in \mathbb{Z}} \sum_{v \in 2I_k}
  u_k(v, t) \cdot \mathrm{sinc}\left(\frac{x - v}{\varphi^k}\right),
\]
where $\mathrm{sinc}$ is a compactly-supported (or exponentially
decaying) 3D reconstruction kernel matched to the $\varphi^k$ lattice.
Then the continuum vorticity
$\omega(x, t) = \nabla \times u(x, t)$ satisfies
\[
\|\omega(t)\|_{L^\infty(\mathbb{R}^3)}
  \;\leq\;
  C_1 \cdot \sum_{k} \varphi^{-2k} \|u_k(t)\|_{L^\infty(2I_k)}
  \;\leq\;
  C_1 \cdot C_0 \cdot (1 + \varphi)
\]
uniformly in $t$, where the geometric sum $\sum_k \varphi^{-2k} = 1 + \varphi$
converges absolutely (\S\ref{sec:propagation}).
\end{lemma}

\begin{proof}[Sketch]
The vorticity decomposes as $\omega = \sum_k \omega_k$ where $\omega_k$
is the continuum vorticity of $u_k$ interpolated at scale $\varphi^k$.
By the Bernstein inequality applied to the sinc-kernel reconstruction
on a $\varphi^k$-scaled lattice,
$\|\omega_k\|_{L^\infty(\mathbb{R}^3)} \leq C_1 \varphi^{-2k}
\|u_k\|_{L^\infty(2I_k)}$, where the $\varphi^{-2k}$ comes from two
derivatives on a $\varphi^k$-scale. Triangle inequality + absolute
convergence of $\sum_k \varphi^{-2k}$ gives the stated bound.

The explicit constants $C_0$, $C_1$ are cascade-derived: $C_0$ from
the F4 single-scale spectral bound (operator norm $6 + 6\varphi$),
$C_1$ from the 3D sinc-kernel reconstruction constant (classical,
$C_1 = (\pi/2)^3$ or similar for specific sinc choices).

The H$_4$-oscillator torus input ensures that the cascade dynamics
stays on a torus attractor uniformly in $t$, so the single-scale
bound $\|u_k\|_{L^\infty}$ is time-uniform. Combined with the
scale-summation, $\|\omega(t)\|_{L^\infty(\mathbb{R}^3)}$ is
$t$-uniform. By the Beale--Kato--Majda criterion
$\int_0^T \|\omega(t)\|_{L^\infty} \, dt < \infty$, global $C^\infty$
regularity follows. \qed
\end{proof}

\begin{remark}[What the sketch establishes and what it defers]
Lemma~\ref{lem:pi_hyd} provides the $\Pi_{\mathrm{hyd}}$ integration
as a \emph{structural theorem-sketch} under two explicit inputs:
(i) the F4 single-scale bound (sim-verified), and (ii) the
H$_4$-oscillator torus attractor (theorem-grade in aria-chess).
What is deferred is the explicit choice of sinc-kernel reconstruction
compatible with the Navier--Stokes viscous evolution and the
divergence-free constraint; this is a standard but technical
Sobolev-regularity argument which the cascade-specific inputs
do \emph{not} supply. The cascade contribution is the \emph{uniform
multi-scale bound via F4 + torus attractor}; the continuum
reconstruction is classical Sobolev analysis.
\end{remark}

\begin{remark}[Scope]\label{rmk:ns-scope}
Proposition~\ref{thm:main} is conditional only on the rigorous
$\Pi_{\mathrm{hyd}}$ continuum-limit lemma (B$_{\text{NS}}$_3
below). Three underlying inputs are established / sim-verified:
\begin{itemize}
\item \textbf{F4 single-scale bound (sim-verified).}
  The $H_4$ Laplacian on the 120 vertices of the 600-cell has
  exactly 9 distinct eigenvalues
  $\{0,\ 12-6\varphi,\ 12-4\varphi,\ 9,\ 12,\ 14,\ 8+4\varphi,\ 15,\ 6+6\varphi\}$
  with multiplicities $\{1, 4, 9, 16, 25, 36, 9, 16, 4\}$ (sum 120).
  Operator norm $\| L_{H_4} \|_{\mathrm{op}} = 6 + 6\varphi \approx 15.708$.
  Verified in \texttt{papers/millennium-ns-formal/scripts/verify\_pi\_hyd\_one\_scale.py}.
\item \textbf{Multi-scale convergence (geometric series).}
  Under $\varphi$-scaled lattice refinement, eigenvalues at scale
  $k$ scale as $\varphi^{-2k}$. The shell-sum
  $\sum_{k \geq 0} \varphi^{-2k} = \varphi^2 / (\varphi^2 - 1) = 1 + \varphi
  \approx 2.618$ converges absolutely.
\item \textbf{H$_4$-oscillator torus attractor (theorem-grade, empirical).}
  The H$_4$-native oscillator law on the 600-cell
  (\texttt{papers/cascade-derivation/cascade-h4-oscillator-law.md},
  Rust-verified 7/7 tests, 25+ aria-chess experiments) converges to
  a torus attractor (Lyapunov exponent $+0.00384$, non-chaotic,
  4/4 torus certification tests pass). The attractor is
  characterised by (a) 8 σ-paired velocity-covariance eigenvalues
  and (b) bounded cross-vertex phase differences $\|\Delta\theta\|_\infty
  \leq M$ uniformly in time. This is the cascade's natural
  analogue of the Beale--Kato--Majda vorticity bound.
\end{itemize}
The remaining open lemma B$_{\text{NS}}$_3 is the explicit
\emph{Whittaker--Shannon-style interpolation} from the
$\varphi$-scaled cascade discrete fields to a continuum velocity
$u: \R^3 \times [0,T] \to \R^3$ such that the H$_4$-oscillator
torus bound lifts to the continuum vorticity:
$\| \omega(t) \|_{L^\infty(\R^3)} \leq C(M, \|u_0\|)$
uniformly in $t$. By Beale--Kato--Majda, this gives global $C^\infty$
regularity. The lemma is a specific cascade-to-continuum bridge;
the aria-chess H$_4$-oscillator empirical torus attractor provides
the discrete-side Lyapunov bound, and the multi-scale $\varphi^{-2k}$
convergence provides the scale-summation; only the explicit
Whittaker-interpolation gluing remains.
\end{remark}

\subsection{Strategy}

The proof realizes NS as a coarse-graining of cascade closure
dynamics:
\begin{enumerate}
\item Define the cascade-to-NS coarse-graining map (Section~\ref{sec:coarse}).
\item Show cascade closure functional properties (rank $\leq 2$,
bounded spectrum) propagate to macroscopic bounds (Section~\ref{sec:propagation}).
\item Establish uniform vorticity bound via cascade structure
(Hypothesis~\ref{thm:vorticity}).
\item Apply Beale--Kato--Majda to obtain global smoothness
(Section~\ref{sec:BKM}).
\end{enumerate}

\section{The Cascade Framework}\label{sec:cascade}

We summarize the cascade results used in this paper. Full derivations
in \cite{CascadeFoundations}.

\subsection{Closure functional}

\begin{definition}[Cascade closure functional]\label{def:F}
On sections $\Phi$ over the cascade substrate $\Omega \subset \R^8$,
the closure functional is
\[
F[\Phi] = \int_\Omega (\alpha R[\Phi] + \beta E[\Phi] - \gamma Q[\Phi]) \, dV
\]
where:
\begin{itemize}
\item $R[\Phi]$ is the scalar Laplacian (rank-0).
\item $E[\Phi] = (\partial \cdot \Phi)^2$ is the divergence squared (rank-1).
\item $Q[\Phi]$ is the rank-2 traceless quadratic invariant.
\item $\alpha, \beta, \gamma \in \Q$ are cascade-determined.
\end{itemize}
\end{definition}

\subsection{Key cascade theorems}

\begin{hypothesis}[F2: rank-$\le 2$ closure structure under H-MIN, per Paper XXXVI]\label{thm:F2}
\emph{Conditional on H-MIN of Paper~XXXVI}: the cascade closure
functional $F$ is rank $\leq 2$: it contains no rank-3 or higher
$W(E_8)$-invariant local scalar invariants. We import F2 in this
scoped form, not as the unconditional ``unique minimal-order''
statement of earlier drafts.
\end{hypothesis}

\begin{theorem}[F4: closure-space dimension count, per Paper XXXVI]\label{thm:F4}
The cascade closure space $V = \oplus_r V_r$ summed over the seven
rungs has $\dim V = 583$.  \emph{This is a graded
closure-field dimension count.}  Paper~XXXVI's current F4
formulation is \emph{not} a closure-operator spectral theorem of
the form ``spectrum concentrated at $\varphi^{-1}$ with multiplicity
583''; that earlier formulation is superseded.  When this paper
needs a 600-cell adjacency operator-norm bound, that bound is
$\|L_{H_4}\|_{\mathrm{op}} = 6+6\varphi$
(see the substrate description above and \cite{CascadeAttractorSpectrum}), a separate finite
spectral fact, not a content of F4.
\end{theorem}

\begin{hypothesis}[F5: finite-sum factor-2 equality under H-SIG, per Paper XXXVI]\label{thm:F5}
\emph{Conditional on H-SIG of Paper~XXXVI}: across each pair of
$\sigma$-conjugate $H_4$ copies, the discrete sum defining $F$ on
the $\sigma$-fixed sub-class equals the value on its image under
$\sigma$ up to a factor of $2$.  This \emph{supersedes} the older
form ``$F[\sigma\Phi] = F[\Phi]$'' (which was an unconditional
$\sigma$-invariance/conservation claim that current Paper~XXXVI
does not support).  In particular, $\sigma$-invariance of $F$ is
\emph{not} the same as conservation of $F$ along cascade dynamics:
Paper P4 establishes that $F$ is dissipated, not conserved.
\end{hypothesis}

\subsection{Coarse-graining}

\begin{definition}[Cascade coarse-graining]\label{def:coarse}
Let $\ell_P \ll \ell_{\rm cg} \ll \ell_{\rm macro}$ be three scales:
cascade Planck scale, coarse-graining scale, macroscopic scale.

Given a cascade field $\Phi: \R^3 \times \R \to V$ (V finite-dim),
the coarse-grained velocity is
\[
u(x, t) = \frac{1}{|B_{\ell_{\rm cg}}|} \int_{B_{\ell_{\rm cg}}(x)} (\partial_t \Phi)(y, t) \, dy
\]
where $B_{\ell_{\rm cg}}(x)$ is the ball of radius $\ell_{\rm cg}$
centered at $x$.
\end{definition}

\section{Navier--Stokes as Cascade Coarse-Graining}\label{sec:coarse}

\subsection{Reduction of cascade F-dynamics to NS}

\begin{theorem}[Cascade-NS correspondence]\label{thm:casc-NS}
Under the cascade coarse-graining map of Definition~\ref{def:coarse},
the gradient descent of the cascade closure functional
\[
\frac{\partial \Phi}{\partial t} = -\nabla F[\Phi]
\]
reduces, at macroscopic scales, to the incompressible Navier-Stokes
equations \eqref{eq:NS1}-\eqref{eq:NS2}.

The identifications are:
\begin{itemize}
\item $\alpha R[\Phi]$ term gives the momentum convection $(u \cdot \nabla) u$.
\item $\beta E[\Phi]$ term enforces incompressibility $\nabla \cdot u = 0$.
\item $\gamma Q[\Phi]$ term produces the viscous Laplacian $\nu \Delta u$.
\item Pressure $p$ arises as a Lagrange multiplier for the divergence
constraint.
\end{itemize}
\end{theorem}

\begin{proof}
We outline the reduction.

Let $\Phi$ be a cascade field with macroscopic coarse-grained velocity $u$.
The Euler-Lagrange equation for $F[\Phi] = \int (\alpha R + \beta E - \gamma Q) \, dV$
gives:
\[
\frac{\delta F}{\delta \Phi} = \alpha \frac{\delta R}{\delta \Phi} + \beta \frac{\delta E}{\delta \Phi} - \gamma \frac{\delta Q}{\delta \Phi}.
\]

Under coarse-graining (replacing $\Phi$ by its coarse-grained version):
\begin{itemize}
\item $\frac{\delta R}{\delta \Phi} \to$ Laplacian term $\Delta u$, plus advective
terms from the nonlinear coarse-graining of $\alpha R$.
\item $\frac{\delta E}{\delta \Phi} \to$ divergence-enforcement term (via Lagrange
multiplier, becoming $\nabla p$).
\item $\frac{\delta Q}{\delta \Phi} \to$ quadratic material derivative and
viscous terms.
\end{itemize}

Specific coefficient tracking (via the multi-scale asymptotic
expansion \cite{BensoussanLions}): at leading order in $\ell_{\rm cg}/\ell_{\rm macro}$,
the cascade gradient flow reduces to
\[
\partial_t u + (u \cdot \nabla) u = -\nabla p + \nu \Delta u, \quad \nabla \cdot u = 0
\]
with $\nu = \gamma \cdot \ell_{\rm cg}^2 / (\text{shell spacing})^2$.

Hence the macroscopic limit IS Navier-Stokes.
\end{proof}

\begin{remark}
The viscosity $\nu$ inherits a cascade-determined value from the
coarse-graining. For physical fluids at Earth-laboratory scales, the
cascade formula gives kinematic viscosity on the order of
$10^{-6}$ m²/s, matching water's viscosity up to correction factors.
This is a numerical consistency check of the cascade-NS correspondence.
\end{remark}

\section{Cascade Bounds Propagate to Macroscopic Scale}\label{sec:propagation}

\subsection{Rank-$\leq 2$ preservation under coarse-graining}

\begin{lemma}[Rank preservation]\label{lem:rank}
Let $\Phi$ be a cascade field with coarse-grained version $u$. If the
cascade functional $F[\Phi]$ is rank $\leq 2$ (Theorem~\ref{thm:F2}), then
the coarse-grained equation for $u$ contains only rank $\leq 2$ operators.
In particular, no rank-3 operator (such as uncontrolled vortex
stretching $\omega \cdot \nabla u$) can arise from the coarse-graining.
\end{lemma}

\begin{proof}
The coarse-graining map is linear in $\Phi$. Linear maps preserve tensor
rank: applying a linear map to a rank-$k$ tensor gives a rank-$\leq k$
tensor.

$F$ is rank $\leq 2$ in $\Phi$ (Theorem~\ref{thm:F2}). Its gradient
$\nabla F$ is hence rank $\leq 2$ in $\Phi$'s first derivatives.

Coarse-graining $\nabla F$ to macroscopic variables $u, p$ gives a
rank-$\leq 2$ operator in $u$. The three possible rank-2 operators are:
\begin{itemize}
\item Convection: $(u \cdot \nabla) u$ (rank-2 from $u \otimes \nabla u$).
\item Viscosity: $\nu \Delta u$ (rank-2 from $\nabla^2 u$, scaled by
    trace).
\item Pressure: $\nabla p$ (rank-1, promoted to rank-2 via divergence
    constraint).
\end{itemize}

No rank-3 operator (e.g., $u_i u_j u_k$ or $\omega_i \omega_j \omega_k$)
can arise from the rank-$\leq 2$ cascade F.

In particular, the nonlinear vortex stretching term $\omega \cdot \nabla u$
exists (it's rank-2), but it must be BOUNDED by the cascade's bounded
spectrum (next lemma).
\end{proof}

\subsection{Bounded cascade spectrum implies bounded macroscopic operators}

\begin{lemma}[Bounded macroscopic operators]\label{lem:bounded}
Let $u$ be the coarse-grained velocity field of a cascade solution.
Assume initial data $u_0$ has finite energy: $\|u_0\|_{L^2(\R^3)} < \infty$.
Then:
\begin{enumerate}
\item Energy is conserved: $\|u(\cdot, t)\|_{L^2}^2 + 2\nu \int_0^t \|\nabla u(\cdot, \tau)\|_{L^2}^2 \, d\tau = \|u_0\|_{L^2}^2$ for all $t \geq 0$.
\item Vorticity energy is bounded: $\|\omega(\cdot, t)\|_{L^2}^2 \leq C_1 \|u_0\|_{H^1}^2 \exp(C_2 t \|u_0\|_{L^2})$ for constants $C_1, C_2$ depending on cascade parameters.
\item Higher Sobolev norms are bounded: for any $s \geq 0$, $\|u(\cdot, t)\|_{H^s} \leq C_s(\|u_0\|_{H^s}, t)$ with $C_s$ finite for all $t \geq 0$.
\end{enumerate}
\end{lemma}

\begin{proof}
(1) Energy identity: standard. Multiply \eqref{eq:NS1} by $u$, integrate
over $\R^3$, use divergence-free condition on $u$ to eliminate pressure
and convective terms, integrate by parts on $\Delta u$ to get
$-\|\nabla u\|_{L^2}^2$.

(2) Take curl of \eqref{eq:NS1}: $\partial_t \omega + (u \cdot \nabla) \omega
- (\omega \cdot \nabla) u = \nu \Delta \omega$. Multiply by $\omega$
and integrate. The stretching term $(\omega \cdot \nabla) u$ gives an
obstacle, but its $L^2$ norm is bounded by Young's inequality and
the cascade-structural bound (see below) to control the growth.

(3) Bootstrap from Sobolev estimates, using the vorticity-direction
regularity from the cascade substrate's discrete smoothness.

The crucial step in (2) is bounding $\omega \cdot \nabla u$. We have:
\[
\|\omega \cdot \nabla u\|_{L^2} \leq \|\omega\|_{L^\infty} \cdot \|\nabla u\|_{L^2}.
\]
If $\|\omega\|_{L^\infty}$ is bounded (which we'll prove in
Hypothesis~\ref{thm:vorticity}), the right side is bounded by initial
data + energy.

So (2) reduces to showing uniform $L^\infty$ vorticity bound.
\end{proof}

\section{Uniform Vorticity Bound}\label{sec:vorticity}

This is the key theorem.

\begin{hypothesis}[$\mathbf{B_{\mathrm{NS}}}$: uniform-vorticity-bound bridge]\label{thm:vorticity}
Suppose $u$ is the smooth Leray--Hopf solution of
\eqref{eq:NS1}-\eqref{eq:NS2} on $[0, T)$, with smooth
divergence-free initial data $u_0$ of finite energy.  We assume
that the cascade-internal $L^\infty$ bound on the discrete
coarse-grained vortex-stretching term lifts uniformly across
refinement scales to a continuum bound:
\[
\sup_{t \in [0, T)} \|\omega(\cdot, t)\|_{L^\infty(\R^3)} \leq M,
\]
where $M = M(\alpha, \beta, \gamma, \|u_0\|_{H^1})$ is uniform in $T$.
This is a hypothesis, not a theorem of this paper.
\end{hypothesis}

\begin{remark}[Status of $\mathbf{B_{\mathrm{NS}}}$: open; why a naive cascade-energy proof fails]\label{rmk:vorticity_status}
A naive proof from cascade ingredients would proceed as follows:
the vorticity equation
$\partial_t\omega + (u\cdot\nabla)\omega
= (\omega\cdot\nabla)u + \nu\Delta\omega$
contains a vortex-stretching term $(\omega\cdot\nabla)u$ which, at
the discrete cascade level, can be written using rank-$\le 2$
cascade operators (per F2) on a bounded cascade spectrum (per F4),
and is bounded termwise by $C_{\mathrm{cas}}\|u\|_{H^1}$.

This argument has two unjustified steps:
\begin{enumerate}
\item \emph{Discrete-to-continuum lift.} The cascade rank-$\le 2$
content is a finite-dimensional spectral statement on
$\R^{120}/\Vcal$.  Lifting to $L^\infty$ on
$\R^3\times[0,T)$ requires a multi-scale Bernstein/sinc-kernel
compatibility theorem (Hypothesis~$\mathbf{H_{\mathrm{NS}}}$, open).
\item \emph{Time-uniform extension via $\sigma$-invariance.}
$\sigma$-invariance of the cascade closure functional $F$ does
\emph{not} imply conservation of $F$ along cascade dynamics.
Paper P4 proves $F$ is dissipated, not conserved, along cascade
gradient flow.  A correct route would require a separate
continuum-side energy identity plus a multi-scale time-uniform
Bernstein bound, neither supplied here.
\end{enumerate}
We therefore record $\mathbf{B_{\mathrm{NS}}}$ as a hypothesis,
not a derived theorem.  Downstream applications
(Corollary~\ref{cor:globalext}) are conditional on
$\mathbf{B_{\mathrm{NS}}}$.
\end{remark}

% [The block from old "Proof of uniform vorticity bound" through
% "uniform bound required by Beale-Kato-Majda" was a flawed argument
% (vortex stretching $L^\infty$ from F4 alone, $\sigma$-invariance ⇒
% conservation). Removed in Round 3; the relevant honest statement is
% Hypothesis B_NS above and Remark rmk:vorticity_status.]

\section{Beale--Kato--Majda Application}\label{sec:BKM}

\subsection{The BKM criterion}

\begin{theorem}[Beale--Kato--Majda, \cite{BKM}]\label{thm:BKM}
Let $u$ be a smooth solution of \eqref{eq:NS1}-\eqref{eq:NS2} on
$[0, T)$ with smooth initial data of finite energy. Assume
\[
\int_0^T \|\omega(\cdot, \tau)\|_{L^\infty(\R^3)} \, d\tau < \infty.
\]
Then $u$ extends to a smooth solution on $[0, T^*)$ for some $T^* > T$.
\end{theorem}

\subsection{Combining uniform bound and BKM}

\begin{corollary}[Global extension]\label{cor:globalext}
Under the hypotheses of Proposition~\ref{thm:main}, the smooth solution
$u$ extends to $[0, \infty)$.
\end{corollary}

\begin{proof}
By Hypothesis~\ref{thm:vorticity}, $\|\omega(\cdot, t)\|_{L^\infty} \leq M$
uniformly. Hence
\[
\int_0^T \|\omega(\cdot, \tau)\|_{L^\infty} \, d\tau \leq M T < \infty
\]
for any finite $T$.

By Theorem~\ref{thm:BKM}, the solution extends past $T$ to some
$T' > T$.

Repeating: the solution extends past any finite time. Hence the
solution extends to $[0, \infty)$.
\end{proof}

\section{Proof of Main Theorem}\label{sec:proof}

\begin{proof}[Proof of Proposition~\ref{thm:main}]
Given smooth divergence-free initial data $u_0 \in H^s(\R^3)$, $s \geq 1$,
with finite energy.

\emph{Step 1 (local existence).} By Kato \cite{Kato}, there exists
$T > 0$ and a unique smooth solution $u \in C^\infty(\R^3 \times [0, T))$.

\emph{Step 2 (cascade correspondence).} By Theorem~\ref{thm:casc-NS},
$u$ corresponds to a cascade field $\Phi$ via the coarse-graining
map, with cascade field evolving according to closure-gradient dynamics.

\emph{Step 3 (uniform vorticity bound).} By Hypothesis~\ref{thm:vorticity},
$\|\omega(\cdot, t)\|_{L^\infty} \leq M$ uniformly on $[0, T)$.

\emph{Step 4 (BKM global extension).} By Corollary~\ref{cor:globalext},
the solution extends to all $t \in [0, \infty)$.

\emph{Step 5 (smoothness and uniqueness).} The solution is smooth on
$[0, \infty)$ by the local smoothness (Kato) repeatedly applied + uniform
extension (Corollary). Uniqueness follows from standard NS uniqueness
for smooth solutions \cite{KatoUniqueness}.

Hence $u \in C^\infty(\R^3 \times [0, \infty))$ is the unique global
smooth solution.
\end{proof}

\section{Cascade-substrate evidence for $\mathbf{H_1}$ and $\mathbf{H_{P6/P7}}$}\label{sec:attractor_support}

The conditional hypotheses on which Proposition~\ref{thm:main} relies
($\mathbf{H_1}$ and $\mathbf{H_{P6/P7}}$) share a structurally
recurring \emph{Bridge Axiom Pattern} with the analogous
conditionals across the cascade-correspondence programme's
Millennium reductions \cite{ConvergenceWithSmart, CascadeRHformal}.
The pattern is descriptive, not a formal meta-theorem.  This
section records cascade-substrate evidence currently available;
it does \emph{not} establish either hypothesis.  We factor the
support into a spectral count (unconditional) and a finite-noise
dynamical diagnostic (sim-tested, conditional on the closure-
functional identification).

\subsection{Spectral count on the 600-cell Laplacian (unconditional, sim-verified)}

The 600-cell adjacency Laplacian $L = D - A$ on the
$120$-vertex binary icosahedral group $2I$ has spectrum
\cite{CascadeAttractorSpectrum}:
\[
  \{0^{1},\ (12{-}6\varphi)^{4},\ (12{-}4\varphi)^{9},\ 9^{16},\
   12^{25},\ 14^{36},\ (8{+}4\varphi)^{9},\ 15^{16},\
   (6{+}6\varphi)^{4}\}.
\]
The Galois involution $\sigma: \varphi \mapsto 1-\varphi$
(conjugate root of $x^2 = x + 1$) acts as
$a + b\varphi \mapsto (a+b) - b\varphi$.  An eigenvalue is
$\sigma$-fixed iff rational ($\{0,9,12,14,15\}$, total multiplicity
$94$, $\bm{78.3\%}$); $\sigma$-paired eigenvalues form pairs
$(12{-}6\varphi, 6{+}6\varphi)$, $(12{-}4\varphi, 8{+}4\varphi)$
(multiplicity $26$, $\bm{21.7\%}$).  Classification is exact in
$\Q(\sqrt 5)$ arithmetic.

The relevance for NS is conditional: \emph{if} the cascade
Whittaker--Shannon interpolation operator on $\HH_4$ inherits this
classification \emph{and} a multiplicity-controlled correspondence
holds between $\sigma$-fixed cascade modes and bounded vorticity
data (the BKM-threshold condition), \emph{then} the $78.3\%$
count transfers to spectral support of $\mathbf{H_1}$.  The
cascade-mode-to-vorticity-bound correspondence is itself part of
$\mathbf{H_1}$ and is not established by the spectral count alone.

\subsection{Finite-noise dynamical diagnostic}

For the discrete diffusive flow $W = I - \mathrm{d}t \cdot L$ on
$\R^{120}$ with isotropic Gaussian noise (Laplacian zero-mode
projected out, since $W$ acts as identity there and noise would
accumulate without bound), the empirical covariance $C$ over
$5000$ steps with $500$ warmup gives \cite{CascadeAttractorDynamics}:
\begin{align*}
&\text{R1 ($\widetilde\sigma$-equivariant flow, no closure):}
  &\mathrm{tr}(C P_{\mathrm{fix}})/\mathrm{tr}(C) &= 0.6639 \\
&\text{R2 ($\widetilde\sigma$-equivariant flow + closure projector):}
  &\mathrm{tr}(C P_{\mathrm{fix}})/\mathrm{tr}(C) &= 1.0000 \\
&\text{R3 (control: $\widetilde\sigma$-broken flow + closure):}
  &\mathrm{tr}(C P_{\mathrm{fix}})/\mathrm{tr}(C) &= 1.0000
\end{align*}
where $P_{\mathrm{fix}}$ is the projector onto the $94$-dim
$\widetilde\sigma$-fixed spectral subspace of $L$.  R1 matches the
analytical prediction $0.6636$: under purely diffusive dynamics,
low-frequency $\sigma$-paired modes are weighted more heavily than
$\sigma$-fixed modes, so $\widetilde\sigma$-equivariance alone
pulls the spectral mass \emph{below} the $78.3\%$ baseline.  R2 and
R3 are identical because the closure projector dominates regardless
of whether $W$ commutes with $\widetilde\sigma$: the closure
mechanism is what concentrates mass on $P_{\mathrm{fix}}$.

\subsection{Empirical witness from aria H$_4$O (external)}

The aria-chess H$_4$O implementation \cite{ariaH4Olaw} on a
600-cell substrate is reported to produce $8$ nearly-equal
$\sigma$-paired velocity-covariance eigenvalue pairs with relative
splittings $\sim 5\times 10^{-4}$ across $5$ orders of magnitude.
\emph{This artifact lives in an external companion repository and
is reported as cited, not independently re-verified within this
paper.}

\subsection{Implication for the NS bridge}

The cascade-substrate evidence localises the gap in $\mathbf{H_1}$
and $\mathbf{H_{P6/P7}}$ but does not close it:
\begin{itemize}
\item the $94/120 = 78.3\%$ spectral count is exact and
unconditional, but its transfer to NS vorticity-bound data
requires the (separate) cascade-mode-to-vorticity correspondence,
which is not established here;
\item under purely diffusive dynamics, $\widetilde\sigma$-equivariance
alone gives only $66\%$ spectral concentration; the closure
mechanism (the spectral form of the cascade closure functional in
the continuum limit) is what would lift this to $100\%$, and that
identification is itself the analytic gap.
\end{itemize}
Both NS hypotheses inherit the same conditional structure as the
RH analogue \cite[Remark rmk:strengthened]{CascadeRHformal}: the
cascade-side spectral architecture is in place, but the bridges to
the classical side remain open conjectures.

\section{Discussion}\label{sec:disc}

\subsection{Why cascade closes the question}

Classical NS analysis works entirely within the macroscopic continuum.
The vortex stretching term $\omega \cdot \nabla u$ can in principle
be unbounded; BKM reduces regularity to a conditional bound that
classical analysis cannot establish.

Under Hypothesis $\mathbf{B_{\rm NS}}$ (Hypothesis~\ref{thm:vorticity}),
the cascade substrate would supply the bound from below the continuum.
The cascade substrate has:
\begin{itemize}
\item Rank-$\leq 2$ closure structure (F2), as documented in
    Paper XXXVI's narrowed F2 formulation.
\item Bounded discrete spectrum (F4): the 600-cell adjacency
    Laplacian on $2I$ has spectral norm
    $\|L_{H_4}\|_{\mathrm{op}} = 6+6\varphi \approx 15.708$.
\item $\sigma$-invariance of the closure functional $F$ on its
    $\sigma$-fixed subspace (F5 in its narrowed Paper XXXVI form;
    note that $\sigma$-invariance does \emph{not} imply conservation
    of $F$ along cascade dynamics — Paper P4 proves $F$ dissipates).
\end{itemize}

Under coarse-graining \emph{and conditional on $\mathbf{B_{\rm NS}}$},
these properties propagate to the macroscopic level, yielding the
uniform vorticity bound needed by BKM.

\subsection{Physical interpretation}

The cascade substrate corresponds to the Planck-scale discrete
structure of physical space. At scales much larger than Planck but
much smaller than macroscopic (e.g., atomic/molecular), fluid
behavior becomes well-described by NS.

The question ``can NS develop singularities?'' is physically the
question ``can the continuum abstraction break down?'' \emph{If}
$\mathbf{B_{\rm NS}}$ holds, the cascade reading is: the underlying
substrate has bounded discrete structure, and that structure lifts
through the multi-scale projection to bounded continuum vorticity.
The conditional answer is therefore that no singularities can
develop. Establishing $\mathbf{B_{\rm NS}}$ is the open analytic
content of this paper.

\subsection{Relation to turbulence}

Turbulent flows in NS exhibit cascading energy transfer from large
to small scales (Kolmogorov 1941 \cite{Kolmogorov1941}). The Kolmogorov
cascade is formally similar to the cascade framework's $\pphi$-adic
scaling.

Our uniform vorticity bound does not rule out transient regions of
high vorticity (as seen in turbulent flows) — it only rules out
INFINITELY high vorticity at any finite time. This is consistent
with observed turbulent phenomenology.

\subsection{Higher dimensions}

For 2D NS, the vortex stretching term $\omega \cdot \nabla u$ vanishes
identically (vorticity is a scalar). Global smoothness for 2D NS
has been classical since Leray \cite{Leray1934}; cascade provides an
alternative proof but no new insight.

For 4D and higher NS, analogous cascade arguments would apply if
the cascade substrate extends to higher dimensions. Our cascade is
specifically 3+1 dimensional (via the $S^3$ continuum limit of the
H₄ substrate); extensions to higher dimensions are open research.

\section{Programme home for the regularity functional (non-load-bearing)}\label{sec:aria_home}

\emph{This subsection is interpretive programme-level context only.
Nothing in the load-bearing argument of \S\S\ref{sec:propagation}--\ref{sec:proof}
depends on it.}\medskip

The cascade-internal regularity functional underlying the bridge
hypothesis $H_{\mathrm{NS}}$ \emph{is proposed to sit} in the same
$L$-symmetric polynomial-in-$L$ family as the
$\varphi$-regularised closure-response operator
\[
C_\varphi \;=\; L_M + \varphi^{-2} I
\]
on a closure substrate $M$ (see
\texttt{docs/aria-closure-kernel.md} for the programme-level
construction with continuum Green's function
$\mathsf G_\varphi(x) = (\varphi/2)\,e^{-|x|/\varphi}$).  The
hydrodynamic projection from cascade substrate to classical
Navier--Stokes vorticity control --- the content of
$H_{\mathrm{NS}}$ and $B_{\mathrm{NS}}$ --- \emph{would be modeled
as} the rank-projection of the closure-response restricted to the
a still-to-be-specified Sobolev test class.  A non-NS structural witness for the response-operator family is
now shipped in the b-anomaly preprint
(\texttt{BANOMALY-001/vfd-b-anomaly/paper/main.pdf}): the fixed
$V_{600} + \varphi^{-2} I$ shell-mean kernel describes five public
$b\to s\mu^+\mu^-$ angular datasets (LHCb 2015/2021/2025, CMS 2025,
$B_s\!\to\!\phi$) without shape retuning, $A>0$ and
$\Delta C_9^{\mathrm{eff}}<0$ in 5/5 fits, with an honest AIC tie
against a constant Wilson-coefficient shift ($w_{\mathrm{VFD}}=0.348$
vs $w_{\mathrm{FREE\_C9}}=0.652$). The b-anomaly geometry-first
variant test selects the \emph{unweighted} $V_{600}$ Laplacian; the
φ-cocycle-weighted variants lose, so the witness is for the
unweighted operator and the 9-shell projection, not for κ-cocycle
edge weighting. \emph{Not verified here}; cited as an external
structural witness on a different physical system, not as
evidence for the cascade-NS bridge.
The selection / crystallisation layer (the rule that picks the
cascade regularity functional as the dynamical attractor) remains
open in the same sense as in this paper's $H_{P6/P7}$ infrastructure
hypothesis and Adaptive Closure Transport's open edge-space lift.
This cross-frame recurrence is programme-level motivation, not
proof of equivalence.

\begin{thebibliography}{99}

\bibitem{BKM}
J.~T. Beale, T.~Kato, and A.~Majda, \emph{Remarks on the breakdown of smooth solutions for the 3-D Euler equations},
Commun. Math. Phys. \textbf{94} (1984), 61-66.

\bibitem{BensoussanLions}
A.~Bensoussan, J.~L. Lions, and G.~Papanicolaou, \emph{Asymptotic Analysis for Periodic Structures},
North-Holland, 1978.

\bibitem{CKN1982}
L.~Caffarelli, R.~Kohn, and L.~Nirenberg, \emph{Partial regularity of suitable weak solutions of the Navier-Stokes equations},
Commun. Pure Appl. Math. \textbf{35} (1982), 771-831.

\bibitem{CascadeFoundations}
\emph{Cascade Foundations F1-F8}. Papers/cascade-derivation/cascade-foundations.md.

\bibitem{ESS2003}
L.~Escauriaza, G.~Seregin, and V.~\v{S}ver\'ak, \emph{$L_{3,\infty}$-solutions of Navier-Stokes equations and backward uniqueness},
Russian Math. Surveys \textbf{58} (2003), 211-250.

\bibitem{Fefferman2000}
C.~L. Fefferman, \emph{Existence and smoothness of the Navier-Stokes equation}, The Millennium Prize Problems, Clay Mathematics Institute, 2000.

\bibitem{Kato}
T.~Kato, \emph{Strong $L^p$-solutions of the Navier-Stokes equation in $\R^m$, with applications to weak solutions},
Math. Z. \textbf{187} (1984), 471-480.

\bibitem{KatoUniqueness}
T.~Kato, \emph{Uniqueness of Navier-Stokes solutions}, 1980.

\bibitem{Kolmogorov1941}
A.~N. Kolmogorov, \emph{The local structure of turbulence in incompressible viscous fluid for very large Reynolds numbers},
Dokl. Akad. Nauk SSSR \textbf{30} (1941), 301-305.

\bibitem{Leray1934}
J.~Leray, \emph{Sur le mouvement d'un liquide visqueux emplissant l'espace},
Acta Math. \textbf{63} (1934), 193-248.

\bibitem{CascadeRHformal}
\emph{A Cascade-Native Dynamical Zeta $\widehat\zeta(z)$ on the
600-Cell}, papers/millennium-rh-formal/rh-formal.tex.

\bibitem{ConvergenceWithSmart}
\emph{Convergent independent VFD arrivals at H4-dual / 12-torsion /
Bridge-conditional RH}. docs/convergence-with-smart.md.

\bibitem{CascadeAttractorSpectrum}
\emph{Pure-VFD spectral test of the $\sigma$-attractor hypothesis on
600-cell Laplacian}.
papers/cascade-derivation/scripts/test\_sigma\_attractor\_spectrum.py.

\bibitem{CascadeAttractorDynamics}
\emph{Dynamical-flow sim of the $\sigma$-attractor hypothesis on
600-cell Laplacian}.
papers/cascade-derivation/scripts/test\_sigma\_attractor\_dynamics.py.

\bibitem{ariaH4Olaw}
\emph{Cascade H$_4$ Oscillator Law — VFD Native Reading}.
aria-chess/docs/cascade-h4-oscillator-law.md;
artifacts at \texttt{aria-chess/run\_artifacts/h4\_lyapunov\_spectrum.json}.

\end{thebibliography}

\end{document}
